% Part: sets-functions-relations
% Chapter: sets
% Section: russells-paradox

\documentclass[../../../include/open-logic-section]{subfiles}


\begin{document}

\olfileid[zh]{sfr}{set}{rus}
\olsection{\zhFullName{Bertrand Russell}{伯特兰·罗素}{罗素}悖论}

外延性容许我们使用记号 $\Setabs{x}{\phi(x)}$，它表示那些满足~$\phi(x)$ 的 $x$ 所构成的\emph{那个}集合。然而，外延性所\emph{真正}容许的只是下述想法。\emph{如果}存在一个集合，其成员恰好是那些 $\phi$，\emph{那么}这样的集合只有一个。换言之：给定某个~$\phi$ 后，集合 $\Setabs{x}{\phi(x)}$ 是唯一的，\emph{如果它存在的话}。

但这个条件句很重要！{}关键在于，并非每个性质都适用于\emph{概括}。也就是说，有些性质\emph{不}定义集合。如果它们全都定义集合，那么我们就会陷入彻头彻尾的矛盾之中。最著名的例子便是\zhFullName{Bertrand Russell}{伯特兰·罗素}{罗素}悖论。

集合可以是其他集合的!!{element}——例如，集合~$A$ 的幂集就是由集合构成的。因此，追问或考察一个集合是否为另一个集合的!!a{element}是有意义的。集合能是其自身的成员吗？关于集合的概念似乎并没有排除这一点。例如，如果\emph{所有}集合构成一个对象的集体，人们也许会认为它们可以被收集成一个集合——所有集合的集合。而它作为一个集合，就会是所有集合之集合的!!a{element}。

当我们考虑「不以其自身为!!a{element}」这一性质，即\emph{非自属的}这一性质时，\zhFullName{Bertrand Russell}{伯特兰·罗素}{罗素}悖论就出现了。如果设想存在一个集合，其中的元素都是不以其自身为!!a{element}的集合，那么会怎样呢？那
\[
R = \Setabs{x}{x \notin x}
\]
存在吗？结果证明，我们可以证明它并不存在。

\begin{thm}[\zhFullName{Bertrand Russell}{伯特兰·罗素}{罗素}悖论]\ollabel{thm:russells-paradox}
	不存在集合 $R = \Setabs{x}{x \notin x}$。
\end{thm}

\begin{proof}
如果 $R = \Setabs{x}{x \notin x}$ 存在，那么 $R \in R$ 当且仅当 $R \notin R$，而这是个矛盾。
\end{proof}

\begin{tagblock}{novice}
\begin{explain}
让我们更缓慢地过一遍这个证明。如果 $R$ 存在，那么追问 $R \in R$ 与否便是有意义的。假设确实有 $R \in R$。现在，$R$~被定义为所有不以其自身为!!{element}的集合所构成的集合。因此，如果 $R \in R$，那么 $R$ 自身并不具有 $R$ 的定义性质。但只有具有该性质的集合才属于~$R$，因此，$R$ 不可能是~$R$ 的!!a{element}，即 $R \notin R$。但 $R$ 不能既是~$R$ 的!!a{element}又不是~$R$ 的元素，于是我们得到一个矛盾。

由于 $R \in R$ 的假设导致矛盾，我们便有 $R \notin R$。但这同样会导致矛盾！{}因为如果 $R \notin R$，那么 $R$ 自身确实具有 $R$ 的定义性质，于是 $R$ 就会像其他所有非自属的集合一样，成为 $R$ 的!!a{element}。而同样地，它不能既不是~$R$ 的!!a{element}、又是~$R$ 的元素。
\end{explain}
\end{tagblock}

\begin{digress}
我们该如何建立一种能避免陷入\zhFullName{Bertrand Russell}{伯特兰·罗素}{罗素}悖论的集合论，即避免作出\emph{不一致的}断言——说 $R = \Setabs{x}{x \notin x}$ 存在——呢？我们需要立下公理，它们给出用于陈述集合何时存在（以及何时不存在）的非常精确的条件。

本章中勾勒的集合论并不这样做。它是\emph{真正朴素的}。它只告诉你，集合服从外延性，并且，如果你有一些集合，那么你可以形成它们的并、交等等。以此更严格地发展集合论是可能的。\oliflabeldef{cumul:::part}{这种严格性将留给 \olref[cumul][][]{part} 这一部分。现在，我们将朴素地进行下去，并小心翼翼地试图避开矛盾。}{}
\end{digress}


\end{document}
